\section{Property Definitions and Parameter Reporting}\label{sec:properties}
\subsection{Species and amount bases}
% P:appendix-properties-01
\Cref{tab:species-basis} distinguishes apparent components from the true species used in equilibrium and transport.
The material, mass, and charge constraints in \cref{eq:conserved-totals,eq:complete-invariants} fix the conversion between these representations.
Molar density and enthalpy are evaluated on the true-species amount basis before conversion to column flow quantities.

\begin{table}[pos=htbp]
\centering\small
\caption{Component and species definitions in the absorber formulation.}
\label{tab:species-basis}
\begin{tabularx}{\linewidth}{@{}p{0.24\linewidth}X@{}}
\toprule
Representation & Components or species\\
\midrule
Apparent liquid & CO$_2$, MEA, H$_2$O\\
True liquid & CO$_2$, MEA, H$_2$O, MEAH$^+$, MEACOO$^-$, HCO$_3^-$, CO$_3^{2-}$, H$_3$O$^+$, OH$^-$\\
Gas & CO$_2$, H$_2$O, N$_2$, O$_2$\\
Interphase transfer & CO$_2$ and H$_2$O\\
Nontransferring components & MEA in the liquid; N$_2$ and O$_2$ in the gas\\
\bottomrule
\end{tabularx}
\end{table}

\subsection{Equilibrium and reference properties}\label{sec:appendix-properties}
% P:appendix-properties-02
Reaction constants are dimensionless and paired with their stated activity conventions.
Temperature correlations retain the coefficient units, logarithm base, and temperature domain of their source representation.
The reference enthalpy and heat-capacity conventions are common to the equilibrium and caloric calculations.
Species identity and order are the same in the reaction system, thermodynamic parameters, reference properties, and film mobility.
% CHECKLIST-BEGIN: method-11
% P:appendix-properties-03
The coefficients in \cref{tab:reactive-constants} use the aqueous infinite-dilution molality convention with water solvent, $m^\circ=\qty{1}{mol\,kg^{-1}}$, and $P^\circ=\qty{1}{bar}$.
For R1--R3 the additional linear-in-temperature coefficient is zero.
R1 and R3 retain additive standard-state shifts $2s$ and $s$, respectively, with $s=4.0165349923$.
The fitted R2 coefficient $a_2$ already includes this conversion, so its additional offset is zero.
The selected R2, R4, and R5 temperature correlations replace their source correlations as complete coefficient sets.

\begin{table}[pos=htbp]
\centering
\caption{Reaction-constant coefficients for \cref{eq:reaction-constants}. Temperature is evaluated in kelvin. R1--R3 include the standard-state offsets; R4 and R5 use their distinct forms.}
\label{tab:reactive-constants}
\small
\begin{tabular}{@{}lrrrr@{}}
\toprule
Reaction & $a_r$ & $b_r$ & $c_r$ & $s_r$ \\
\midrule
R1 & 132.899 & $-13445.9$ & $-22.4773$ & 8.0330699846 \\
R2 & 232.331415338844 & $-11105.6400305203$ & $-36.7816$ & 0 \\
R3 & 216.049 & $-12431.7$ & $-35.4819$ & 4.0165349923 \\
\midrule
R4 & \multicolumn{4}{l}{$a_4=1.505015374192114$, $b_4=-1317.0489707842564$} \\
R5 & \multicolumn{4}{l}{$A_5=3037.6399534696106$, $B_5=-1.0173150837285996$, $C_5=0.0004277$} \\
\bottomrule
\end{tabular}
\end{table}

\begin{table}[pos=htbp]
\centering\small
\caption{Temperature intervals assigned to the reaction correlations. A model evaluation interval does not enlarge the underlying source interval.}
\label{tab:reaction-temperature-intervals}
\begin{tabular}{@{}ll@{}}
\toprule
Reaction & Evaluation interval (K)\\
\midrule
R1, R3 & 273.15--498.15\\
R2, R4 & 293.15--393.15\\
R5 & 293.15--393.15\\
\bottomrule
\end{tabular}
\end{table}

% P:appendix-properties-04
The underlying source intervals are \qtyrange{273.15}{498.15}{K} for the R1--R3 forms, \qtyrange{293.15}{323.15}{K} for R4, and \qtyrange{273.15}{323.15}{K} for R5.
Use of R4 or R5 above \qty{323.15}{K} extrapolates the respective source relation.
% CHECKLIST-REVIEW: Typed fitted R2/R4/R5 from parameter a9186c93 replace complete source correlations; no second R2 offset. Retain primary-source/domain review and match the final run identity; see SOURCE_MAP.md.

% CHECKLIST-END: method-11
% CHECKLIST-BEGIN: method-10
% P:appendix-properties-05
The ePC-SAFT description uses the nine pure-component records in \cref{tab:reactive-pure-parameters}, general-site association, Debye--H\"uckel electrostatics, the original Born contribution, and solvent-only MEA/water relative-permittivity mixing.
For the dispersion combining rule,
\begin{align}
\epsilon_{ij}=\sqrt{\epsilon_i\epsilon_j}(1-k_{ij}),\qquad k_{ij}=k_{ji}.
\end{align}
All like-charge ionic pairs have $k_{ij}=1$ and therefore zero unlike-pair dispersion energy.
All other distinct pairs have $k_{ij}=0$ except those listed in \cref{tab:reactive-pair-parameters}; together these rules specify all 36 distinct pairs.
The CO$_2$/MEA reciprocal-temperature slope is zero.

\begin{table}[pos=htbp]
\centering\small
\caption{Nonzero dispersion corrections outside the like-charge rule. Values are dimensionless and symmetric in species indices.}
\label{tab:reactive-pair-parameters}
\begin{tabularx}{\linewidth}{@{}lrX@{}}
\toprule
Pair & $k_{ij}$ & Parameter treatment\\
\midrule
\ce{CO2}/\ce{H2O} & 0.0132628791769 & Fixed mixture correction\\
\ce{MEA}/\ce{H2O} & $-0.0735274987499$ & Fitted to neutral MEA/water VLE\\
\ce{MEAH+}/\ce{MEACOO-} & $-0.00201813457644$ & Fixed historical electrolyte value\\
\ce{CO3^2-}/\ce{H2O} & $-0.25$ & Assigned ion/water correction\\
\ce{H3O+}/\ce{H2O} & 0.25 & Assigned ion/water correction\\
\ce{OH-}/\ce{H2O} & $-0.25$ & Assigned ion/water correction\\
\bottomrule
\end{tabularx}
\end{table}

% P:appendix-parameter-origin
The MEA/water binary correction is fitted to a 25-point neutral-mixture VLE dataset using pressure-level groups for model selection and a final fit to the complete dataset.
The other nonzero pair values in \cref{tab:reactive-pair-parameters} are fixed inputs: the molecular CO$_2$/water correction and ion/water corrections are specified mixture assumptions, and the MEAH$^+$/MEACOO$^-$ correction is retained from the starting electrolyte parameterization.
Like-charge dispersion exclusion follows the selected combining rule \cite{Figiel2025}.
These distinct origins separate parameter estimation from the physical assumptions used in the column calculation.

% P:appendix-parameter-fit
The CO$_2$ dispersion energy is selected against pressure and speciation measurements, while the remaining pure-component segment and ionic diameter values are adopted from the starting parameterization.
The R2, R4, and R5 temperature coefficients are adjusted through reaction-enthalpy shifts while preserving their equilibrium constants at \qty{313.15}{K}.
The reaction-temperature fit uses pressure and speciation targets together with six calorimetric intervals from Kim and Svendsen \cite{Kim2007}; its pressure targets include Jou et al. measurements at \qty{393.15}{K} \cite{Jou1995}.
The Xu pressure measurements and the calorimetry block at \qty{393.15}{K} are excluded from the fitting objective and used for separate thermodynamic evaluation \cite{Xu2011,Kim2007}, while the Kim et al. measurements provide a model-selection comparison \cite{Kim2014}.
The resulting reaction functions and the neutral caloric anchors jointly determine the ionic reference enthalpies and heat capacities.
Absorber capture and temperature profiles provide a separate column evaluation rather than the targets of these property fits.
% CHECKLIST-REVIEW: Parameter a9186c93 source IDs cai-1996-neutral-refit, mea-best-in-slot-campaign-2026-09-02, and mea-reaction-temperature-fit-2026-09-03 identify these fits. Preserve upstream fit/source files and complete primary bibliography locators in the final data deposit.

% P:appendix-properties-06
MEA, water, and CO$_2$ each have one donor-labelled site $a$ and one acceptor-labelled site $b$ in the selected topology.
The explicit association pairs are reported in \cref{tab:reactive-association}; site labels denote the model topology, including induced association of CO$_2$.
The two MEA/water cross pairs use arithmetic averaging of association energies and geometric averaging of association volumes.
No other site pairs are active, and the ionic species have no association sites in this parameterization.

\begin{table}[pos=htbp]
\centering\small
\caption{Explicit association parameters. MEA/water cross association uses the combining rules stated in the text.}
\label{tab:reactive-association}
\begin{tabular}{@{}lrr@{}}
\toprule
Site pair & $\epsilon^{AB}/k$ (K) & $\kappa^{AB}$\\
\midrule
MEA $a$ / MEA $b$ & 2586.3 & 0.03747\\
Water $a$ / water $b$ & 2425.7 & 0.04509\\
CO$_2$ $a$ / water $b$ & 1212.85 & 0.04509\\
CO$_2$ $b$ / water $a$ & 1212.85 & 0.04509\\
\bottomrule
\end{tabular}
\end{table}

% P:appendix-properties-07
The MEA relative permittivity is 32, and the ionic-region value is 8.
For water, with $\Delta T=T-298.15\,\mathrm K$ evaluated numerically in kelvin,
\begin{align}
\varepsilon_{r,W}=78.4104347858-0.361337662337\Delta T
+0.00076555618295(\Delta T)^2.
\end{align}
The salt-free solvent relative permittivity uses mass fractions normalized over the neutral MEA and water solvent species \cite{Figiel2025}:
\begin{align}
w_i^{\mathrm{solv}}&=\frac{n_iM_i}{n_A M_A+n_W M_W},\qquad i\in\{A,W\},\\
\varepsilon_r(T,\mathbf n)&=32w_A^{\mathrm{solv}}+\varepsilon_{r,W}(T)w_W^{\mathrm{solv}}.
\label{eq:solvent-permittivity}
\end{align}
Here $n_A$ is the amount of neutral MEA, rather than the total apparent amine amount.
CO$_2$ and ionic species are excluded from this solvent normalization.
The common parameter-evaluation interval is \qtyrange{293.15}{393.15}{K}; the individual reaction-source intervals are distinguished in \cref{tab:reaction-temperature-intervals}.
% CHECKLIST-REMAINING: Add the selected eNRTL pure/pair and model-specific reference records. The ePC-SAFT coefficient treatment and solvent-only dielectric equation are documented; complete primary-source locators and match final run identities in SOURCE_MAP.md.

\begin{table}[pos=htbp]
\centering\small
\caption{Additional thermodynamic inputs for the activity-coefficient comparison and solvent dielectric mixing. The ePC-SAFT pair and association coefficients are given in the preceding tables.}
\label{tab:comparison-thermo-inputs}
\begin{tabularx}{\linewidth}{@{}lX@{}}
\toprule
Input & Definition and coefficients\\
\midrule
eNRTL species and reactions & CO$_2$, MEA, H$_2$O, MEAH$^+$, MEACOO$^-$, HCO$_3^-$; the two formation reactions in \cref{eq:enrtl-combined-reactions}\\
eNRTL interactions & \textit{[Molecule/molecule, molecule/ion-pair, and ion-pair records]}\\
eNRTL temperature dependence & \textit{[Interaction forms, nonrandomness, and coefficient units]}\\
Solvent dielectric mixing & Neutral-solvent mass fractions; \cref{eq:solvent-permittivity}\\
eNRTL reference conversion & \textit{[Model-specific solute reference offsets]}\\
eNRTL parameter sources & \textit{[Source records and applicable domains]}\\
\bottomrule
\end{tabularx}
\end{table}

% CHECKLIST-END: method-10
% INSERT Table: reference enthalpies/heat capacities, units and amount basis, standard pressure, temperature ranges, sources, and the fixed-P/conserved-total equilibrium derivative convention.
% Preserve any source-domain extrapolation explicitly alongside the affected parameter; do not hide it in workflow prose.

\subsection{Transport and packing correlations}
% P:appendix-properties-08
The Henry-law reference uses a concentration-based coefficient consistent with \cref{eq:henry}; the N$_2$O analogy provides the solvent-mixture solubility relation \cite{Jiru2012}.
Viscosity, surface tension, and diffusivity correlations supply the physical transport properties \cite{Morgan2015}.
Packing area, void fraction, hydraulic diameter, and correlation coefficients define the hydraulic model \cite{Billet1999,Tsai2010,Chinen2018}.
Within each controlled comparison, these inputs remain common except for the property group explicitly varied in \cref{tab:comparison-design}.

% P:appendix-properties-09
The mobility diffusivities $D_i$ describe individual true species, whereas $D_C^{\mathrm{column}}$ defines the physical film thickness in \cref{eq:film-thickness}.
The distinction preserves the respective meanings of the mobility approximation and the packed-column mass-transfer correlation.
% CHECKLIST-BEGIN: method-13
% P:appendix-properties-10
The coefficient tables below define the common-property calculation on an apparent liquid-component basis and a four-component gas basis.
Temperature is in kelvin unless a Celsius variable is explicitly defined, and empirical coefficients absorb the units needed by their stated expressions.
The selected holdup relation is
\begin{align}
h_l^*&=11.4474\left[3.185966\,u_l(\mu_l/\rho_{l,m})^{1/3}\right]^{0.6471},\\
h_l&=\frac{\varepsilon h_l^*}{\varepsilon+h_l^*},\qquad h_v=\varepsilon-h_l,
\label{eq:common-holdup}
\end{align}
using $u_l$ in \unit{m\,s^{-1}}, $\mu_l$ in \unit{Pa\,s}, and $\rho_{l,m}$ in \unit{kg\,m^{-3}}.
The rational mapping keeps the liquid holdup below the packing void fraction.
For gas component $j$, the fugacity-based transfer coefficient and the heat-transfer analogy are
\begin{align}
k_{g,j}&=\frac{C_V}{RT_v}\sqrt{\frac{a_p}{d_hh_v}}D_{v,j}^{2/3}
\left(\frac{\mu_v}{\rho_{v,m}}\right)^{1/3}
\left(\frac{u_v\rho_{v,m}}{a_p\mu_v}\right)^{3/4},\\
U_T&=\left[\frac{(k_{g,C}P)^3 k_{t,v}^{\,2}c_{p,v}}
{(\rho_vD_{v,C})^2}\right]^{1/3}.
\label{eq:common-gas-heat-transfer}
\end{align}
Here $\rho_v$ is gas molar density, $\rho_{v,m}$ is gas mass density, $c_{p,v}$ is molar heat capacity, and $k_{t,v}$ is gas thermal conductivity.
The resulting units are \unit{mol\,m^{-2}\,s^{-1}\,Pa^{-1}} for $k_{g,j}$ and \unit{W\,m^{-2}\,K^{-1}} for $U_T$.
% CHECKLIST-REVIEW: Complete individual coefficient-source locators and validity domains, and reconcile these common-property definitions with the final comparison settings.
% CHECKLIST-END: method-13
% INSERT final model-specific property equations where needed to reproduce the reported comparisons.

\subsection{Henry-law solvent-mixture correlation}\label{app:henry}
% P:appendix-properties-11
The concentration-based N$_2$O analogy retains separate pure-solvent relations and a mixture correction \cite{Jiru2012}.
With $T$ in kelvin, the coefficients below give $H^c$ in \unit{Pa\,m^3\,mol^{-1}}:
\begin{align}
H_{\mathrm{N_2O,MEA}}^c&=2.448\times10^5\exp(-1348/T),\\
H_{\mathrm{CO_2,H_2O}}^c&=3.52\times10^6\exp(-2113/T),\\
H_{\mathrm{N_2O,H_2O}}^c&=8.449\times10^6\exp(-2283/T),\\
H_{\mathrm{CO_2,MEA}}^c&=H_{\mathrm{N_2O,MEA}}^c
\frac{H_{\mathrm{CO_2,H_2O}}^c}{H_{\mathrm{N_2O,H_2O}}^c}.
\end{align}
Let $w_A$ and $w_W$ be MEA and water mass fractions normalized over the CO$_2$-free solvent.
The solvent-mixture relation is
\begin{align}
\ln\left(\frac{H_C^c}{H_*}\right)
&=w_A\ln\left(\frac{H_{\mathrm{CO_2,MEA}}^c}{H_*}\right)
+w_W\ln\left(\frac{H_{\mathrm{CO_2,H_2O}}^c}{H_*}\right)+R_a,\\
R_a&=\left[a_1+a_2(T-273.15)+a_3(T-273.15)^2+bw_W\right]w_Aw_W,
\end{align}
where $H_*=\qty{1}{Pa\,m^3\,mol^{-1}}$ makes the logarithm arguments dimensionless.
The mixture coefficients are $a_1=1.70981$, $a_2=0.03972$, $a_3=-4.3\times10^{-4}$, and $b=-2.20377$, with temperature units absorbed in $a_2$ and $a_3$.
This relation belongs to the Henry reference; the ePC-SAFT arm obtains fugacity from the equation of state.

\subsection{Common-property density correlation}
% P:appendix-properties-12
The common-property comparison uses an apparent-component liquid-volume correlation \cite{Morgan2015}.
Writing $x_C^{\mathrm{app}}$, $x_A^{\mathrm{app}}$, and $x_W^{\mathrm{app}}$ for apparent mole fractions, the volume contributions are
\begin{align}
\widetilde V_C&=a_C+(b_C+c_Cx_A^{\mathrm{app}})x_A^{\mathrm{app}},\\
\widetilde V_A&=\frac{1000M_A}{a_AT^2+b_AT+c_A}
+x_W^{\mathrm{app}}(d_A+e_Ax_A^{\mathrm{app}}),\\
\widetilde V_W&=\frac{1000M_W}{a_WT^2+b_WT+c_W},\\
V^{\mathrm{app}}&=10^{-6}\sum_{j\in\{C,A,W\}}x_j^{\mathrm{app}}\widetilde V_j.
\end{align}
The $\widetilde V_j$ are in \unit{mL\,mol^{-1}}, $M_j$ are in \unit{kg\,mol^{-1}}, and $V^{\mathrm{app}}$ is in \unit{m^3\,mol^{-1}}.
The associated densities are
\begin{align}
\rho_{\mathrm{mol}}^{\mathrm{app}}=\frac1{V^{\mathrm{app}}},\qquad
\rho_m=\frac{\sum_jx_j^{\mathrm{app}}M_j}{V^{\mathrm{app}}}.
\end{align}
This apparent molar density is used on its defined component basis; it is distinct from the true-species density in \cref{eq:true-concentration}.
The full-property ePC-SAFT comparison uses the equation-of-state volume consistently with its equilibrium species amounts.

\begin{table}[pos=htbp]
\centering\small
\caption{Coefficients of the apparent liquid-volume reference correlation. Coefficient units correspond to the temperature and volume conventions in the equations above.}
\begin{tabular}{@{}lrrrrr@{}}
\toprule
Component & $a$ & $b$ & $c$ & $d$ & $e$\\
\midrule
CO$_2$ & 10.2074 & 207.0 & $-563.3701$ & -- & --\\
MEA & $-5.35162\times10^{-7}$ & $-4.51417\times10^{-4}$ & 1.19451 & $-2.2642$ & 3.0059\\
H$_2$O & $-3.2484\times10^{-6}$ & 0.00165 & 0.793 & -- & --\\
\bottomrule
\end{tabular}
\end{table}

\subsection{Surface tension}
% P:appendix-properties-13
The mixture surface tension is expressed through pure-component contributions and loading-dependent correction functions \cite{Morgan2015}:
\begin{align}
\sigma_j&=s_{1,j}(1-T/T_{c,j})^{s_{2,j}+s_{3,j}(T/T_{c,j})+s_{4,j}(T/T_{c,j})^2},\qquad j\in\{A,W\},\\
\sigma_C&=S_1r^2+S_2r+S_3+T(S_4r^2+S_5r+S_6),\\
f_\sigma&=a+b\alpha+c\alpha^2+dr+er^2,\\
g_\sigma&=f+g\alpha+h\alpha^2+ir+jr^2,\\
\sigma_l&=\sigma_W+(\sigma_C-\sigma_W)f_\sigma x_C^{\mathrm{app}}
+(\sigma_A-\sigma_W)g_\sigma x_A^{\mathrm{app}}.
\end{align}
Here $r=w_A/(w_A+w_W)$ and $\alpha$ is the apparent CO$_2$/MEA loading.
The surface tension is in \unit{N\,m^{-1}} and enters the wetting relation in \cref{eq:effective-area}.
The correction functions are empirical mixture terms and are evaluated on the composition basis for which their coefficients are defined.

\begin{table}[pos=htbp]
\centering\small
\caption{Surface-tension mixture coefficients, using $T$ in kelvin and the dimensionless solvent mass fraction $r$.}
\label{tab:surface-mixture-coefficients}
\begin{tabular}{@{}lrlr@{}}
\toprule
Coefficient & Value & Coefficient & Value\\
\midrule
$S_1$ & $-0.00589934906112609$ & $a$ & 1070.65668317975\\
$S_2$ & 0.00175020536428591 & $b$ & $-2578.78134208703$\\
$S_3$ & 0.129650182728177 & $c$ & 3399.24113311222\\
$S_4$ & 0.0000126444768126308 & $d$ & $-2352.47410135319$\\
$S_5$ & $-0.00000573954817199691$ & $e$ & 2960.24753687833\\
$S_6$ & $-0.00018969005534195$ & $f$ & 3.06684894924048\\
 & & $g$ & $-1.79435372759593$\\
 & & $h$ & $-7.2124219075848$\\
 & & $i$ & 2.97502322396621\\
 & & $j$ & $-10.5738529301824$\\
\bottomrule
\end{tabular}
\end{table}

% P:appendix-properties-14
The pure-solvent coefficient sets are
\begin{align}
(s_1,s_2,s_3,s_4,T_c)_A&=(0.09945,1.067,0,0,614.45),\\
(s_1,s_2,s_3,s_4,T_c)_W&=(0.18548,2.717,-3.554,2.047,647.13),
\end{align}
with $T_c$ in kelvin and $s_1$ in \unit{N\,m^{-1}}.

\subsection{Liquid and vapor viscosity}
% P:appendix-properties-15
For the aqueous liquid, the water reference and mixture factor are
\begin{align}
\mu_W&=A_\mu10^{B_\mu[293.15-T-C_\mu(T-293.15)^2]/(T-D_\mu)},\\
\mu_l&=\mu_W\exp\left\{
\frac{r_\%[T(ar_\%+b)+cr_\%+d][\alpha(er_\%+fT+g)+1]}{T^2}
\right\},
\end{align}
with $r_\%=100w_A/(w_A+w_W)$.
The reference coefficients are $A_\mu=\qty{1.002e-3}{Pa\,s}$, $B_\mu=1.3272$, $C_\mu=0.001053$, and $D_\mu=\qty{168.15}{K}$.
The mixture coefficients retain the empirical temperature and solvent-concentration units of this relation \cite{Morgan2015}.
Their values are listed in \cref{tab:viscosity-mixture-coefficients}.

\begin{table}[pos=htbp]
\centering\small
\caption{Liquid-viscosity mixture coefficients for $r_\%$ in mass percent and $T$ in kelvin.}
\label{tab:viscosity-mixture-coefficients}
\begin{tabular}{@{}lr@{}}
\toprule
Coefficient & Value\\
\midrule
$a$ & $-0.0854041877181552$\\
$b$ & 2.72913373574306\\
$c$ & 35.1158892542595\\
$d$ & 1805.52759876533\\
$e$ & 0.00716025669867574\\
$f$ & 0.0106488402285381\\
$g$ & $-0.0854041877181552$\\
\bottomrule
\end{tabular}
\end{table}

% P:appendix-properties-16
For the gas, the mixture relation is
\begin{align}
\mu_v&=\sum_i\frac{y_i\mu_i}{\sum_jy_j\Phi_{ij}},\\
\Phi_{ij}&=\frac{\left[1+\sqrt{\mu_i/\mu_j}(M_j/M_i)^{1/4}\right]^2}
{\sqrt{8(1+M_i/M_j)}},\qquad \Phi_{ii}=1.
\end{align}
All four gas components enter the sums.
The pure-gas viscosity expressions use the common forms
\begin{align}
\mu_i(T)&=\frac{A_iT^{B_i}}{1+C_i/T},
\qquad\text{or}\qquad
\mu_i(T)=\mu_{i,0}\frac{T_{i,0}+S_i}{T+S_i}
\left(\frac{T}{T_{i,0}}\right)^{3/2}.
\end{align}
The correlation form and coefficients are selected by species, with viscosity in \unit{Pa\,s}.
The four pure-gas expressions are
\begin{align}
\mu_C&=\frac{2.148\times10^{-6}T^{0.46}}{1+290/T},&
\mu_W&=1.7096\times10^{-8}T^{1.1146},\\
\mu_N&=1.781\times10^{-5}\frac{300.55+111}{T+111}
\left(\frac{T}{300.55}\right)^{3/2},\\
\mu_O&=2.018\times10^{-5}\frac{292.25+127}{T+127}
\left(\frac{T}{292.25}\right)^{3/2}.
\end{align}

\subsection{Molecular diffusivity}
% P:appendix-properties-17
For the liquid CO$_2$ correlation, define $c_A^{(L)}=10^{-3}c_A$ in \unit{mol\,L^{-1}} when $c_A$ is in \unit{mol\,m^{-3}}.
The diffusivity expressions are
\begin{align}
D_C&=(a+b c_A^{(L)}+c[c_A^{(L)}]^2)
\exp\left(\frac{d+e c_A^{(L)}}{T}\right),\\
D_A&=\exp(a_A+b_A/T+c_A^*c_A),\\
D_{\mathrm{ion}}&=\exp(a_I+b_I/T+c_I\ln\mu_l^*).
\end{align}
Here $\mu_l^*$ is viscosity evaluated numerically in \unit{Pa\,s}, and the coefficients give diffusivity in \unit{m^2\,s^{-1}}.
The empirical common-property ion correlation and the species-resolved mobility inputs have separate definitions.
The coefficient sets, in the order used above, are
\begin{align}
(a,b,c)&=(2.35\times10^{-6},\ 2.9837\times10^{-8},\ -9.7078\times10^{-9}),\\
(d,e)&=(-2119,\ -20.132),\\
(a_A,b_A,c_A^*)&=(-13.275,\ -2198.3,\ -7.8142\times10^{-5}),\\
(a_I,b_I,c_I)&=(-22.64,\ -1000,\ -0.7).
\end{align}

% P:appendix-properties-18
For the gas, the binary and mixture diffusion coefficients are
\begin{align}
D_{ij}&=\frac{1.013\times10^{-2}T^{1.75}}{P}
\frac{\sqrt{10^{-3}(M_i^{-1}+M_j^{-1})}}
{(V_{D,i}^{1/3}+V_{D,j}^{1/3})^2},\\
D_i&=\frac{1-y_i}{\sum_{j\ne i}y_j/D_{ij}}.
\end{align}
These expressions use $P$ in pascals, $M_i$ in \unit{kg\,mol^{-1}}, and the tabulated diffusion-volume convention.
The diffusion volumes for CO$_2$, water, N$_2$, and O$_2$ are 26.7, 13.1, 18.5, and 16.3, respectively.

\subsection{Heat-capacity and thermal-conductivity references}
% P:appendix-properties-19
The common-property sensible heat-capacity correlations use temperature polynomials.
For MEA and water, with $t$ evaluated numerically in degrees Celsius,
\begin{align}
c_{p,j}^{\mathrm{ref}}=1000M_j\sum_{k=0}^{4}a_{jk}t^k,
\qquad j\in\{A,W\}.
\end{align}
The result is molar heat capacity in \unit{J\,mol^{-1}\,K^{-1}} \cite{Hilliard2008}.

\begin{table}[pos=htbp]
\centering\small
\caption{Solvent sensible-heat-capacity coefficients for $t$ in degrees Celsius.}
\label{tab:solvent-cp-coefficients}
\begin{tabular}{@{}lrrrrr@{}}
\toprule
Species & $a_0$ & $a_1$ & $a_2$ & $a_3$ & $a_4$\\
\midrule
MEA & 2.6161 & $3.706\times10^{-3}$ & $3.787\times10^{-6}$ & 0 & 0\\
Water & 4.2107 & $-1.696\times10^{-3}$ & $2.568\times10^{-5}$ & $-1.095\times10^{-7}$ & $3.038\times10^{-10}$\\
\bottomrule
\end{tabular}
\end{table}

% P:appendix-properties-20
For the gas reference,
\begin{align}
c_{p,i}^{\mathrm{ref}}=R(A_i+B_iT+C_iT^{-2}),\qquad
c_{p,v}^{\mathrm{ref}}=\sum_i y_i c_{p,i}^{\mathrm{ref}}.
\end{align}

\begin{table}[pos=htbp]
\centering\small
\caption{Gas sensible-heat-capacity coefficients for $c_p/R=A+BT+CT^{-2}$ with $T$ in kelvin.}
\label{tab:gas-cp-coefficients}
\begin{tabular}{@{}lrrr@{}}
\toprule
Species & $A$ & $B$ (K$^{-1}$) & $C$ (K$^2$)\\
\midrule
CO$_2$ & 5.457 & $1.045\times10^{-3}$ & $-1.157\times10^5$\\
Water & 3.470 & $1.450\times10^{-3}$ & $1.210\times10^4$\\
N$_2$ & 3.280 & $0.593\times10^{-3}$ & $4.000\times10^3$\\
O$_2$ & 3.639 & $0.506\times10^{-3}$ & $-2.270\times10^4$\\
\bottomrule
\end{tabular}
\end{table}

% P:appendix-properties-21
Integration with \cref{eq:reference-enthalpy} gives the corresponding sensible enthalpy change.
These sensible-property relations supply part of the reference comparison.
For all nine reactive species, the complete enthalpy construction takes the form
\begin{align}
H(T,P,\mathbf n)=\sum_{i=1}^{9}n_i h_i^\circ(T)+H^{\mathrm{dep}}(T,P,\mathbf n),
\end{align}
where $H^{\mathrm{dep}}$ contains the phase and mixture contributions relative to the same species reference states.
The species reference enthalpies are constrained by the temperature dependence of the reaction constants after conversion to the EOS reference convention,
\begin{align}
\boldsymbol\nu^{\mathsf T}\mathbf h^\circ(T)
&=\boldsymbol{\Delta h}_{r}^{\mathrm{ref}}(T),&
\Delta h_{r}^{\mathrm{ref}}(T)&=RT^2\odv{\ln K_r^{\mathrm{ref}}}{T}.
\label{eq:reaction-reference-enthalpy}
\end{align}
Here $K_r^{\mathrm{ref}}$ includes the temperature-dependent source-reference transfer and the conversion of the activity basis; differentiating only the unconverted source correlation would omit these terms.
Consequently, reference enthalpies, their heat-capacity derivatives, and reaction equilibrium cannot be assigned independently.
The equilibrium caloric calculation also includes the reactive-composition contribution stated in \cref{subsec:calorics}.
% CHECKLIST-BEGIN: method-12
% P:appendix-properties-22
\Cref{tab:species-reference-calorics} specifies the species reference enthalpies and heat-capacity functions on the common amount and pressure basis.
The nine functions are reconstructed from three neutral anchors, five reaction relations, and one ionic reference convention.
They do not require nine independently measured species heat capacities.
The neutral anchors use ideal-gas CO$_2$ and pure-liquid MEA and water.
For each liquid anchor $j$, the EOS residual at the anchor pressure $P_a$ is removed before the reference function is combined with the mixture residual:
\begin{align}
h_j^\circ(T)&=h_j^{\mathrm{anchor}}(T)-h_j^{\mathrm{res,pure}}(T,P_a),\\
c_{p,j}^\circ(T)&=c_{p,j}^{\mathrm{anchor}}(T)-c_{p,j}^{\mathrm{res,pure}}(T,P_a).
\label{eq:neutral-caloric-anchor}
\end{align}
Thus adding the residual mixture enthalpy recovers the specified pure-liquid anchors without counting their residual contribution twice.
The neutral heat-capacity functions determine the temperature dependence of the anchors, while formation enthalpies fix their integration constants.

% P:appendix-properties-23
Let $\mathbf S$ select CO$_2$, MEA, and water from the nine-species vector, and let $\mathbf g=\mathbf e_{\mathrm{H_3O^+}}-\mathbf e_{\mathrm{H_2O}}$.
At each temperature, the reference functions satisfy
\begin{align}
\begin{bmatrix}\boldsymbol\nu^{\mathsf T}\\\mathbf S\\\mathbf g^{\mathsf T}\end{bmatrix}
\mathbf h^\circ(T)
&=\begin{bmatrix}\boldsymbol{\Delta h}_{r}^{\mathrm{ref}}(T)\\\mathbf h_N^\circ(T)\\0\end{bmatrix},&
\mathbf c_p^\circ(T)&=\odv{\mathbf h^\circ}{T}.
\label{eq:caloric-reconstruction}
\end{align}
The last row sets $h_{\mathrm{H_3O^+}}^\circ=h_{\mathrm{H_2O}}^\circ$ as the ionic reference convention; it is not an additional measured property.
Differentiating the same system gives the ionic heat capacities consistently with the reaction correlations and neutral anchors.
When polynomial functions represent the reconstructed enthalpies, their analytic derivatives define the heat capacities; reconstruction accuracy is checked between the temperature points used to determine the coefficients.
The fixed-composition heat capacity is
\begin{align}
C_{P,\mathbf n}=\sum_i n_i c_{p,i}^\circ(T)+C_P^{\mathrm{dep}}(T,P,\mathbf n).
\end{align}
The equilibrium heat capacity additionally includes the speciation term in \cref{subsec:calorics}.
Gas, liquid, and transferred-material enthalpies use a common reference: any change from formation enthalpies to a sensible-enthalpy zero is applied consistently across phases and preserves the reaction relations.
% CHECKLIST-REVIEW: This v2 reconstruction has fingerprint 6919acbc and matches selected parameter a9186c93; generation Engine 8438ce5f / wheel 40fba7cf. All reaction and caloric coefficients were updated together. Match final run identities; SOURCE_MAP.md retains exact provenance.

The numerical representation uses $T_0=\qty{353.15}{K}$, $P_a=\qty{101325}{Pa}$, and $\theta=(T-T_0)/(100\,\mathrm K)$:
\begin{align}
c_{p,i}^\circ(T)&=\sum_{k=0}^{7}b_{ik}\theta^k,\\
h_i^\circ(T)&=h_i^\circ(T_0)+(100\,\mathrm K)\sum_{k=0}^{7}\frac{b_{ik}}{k+1}\theta^{k+1}.
\label{eq:caloric-polynomial}
\end{align}
The polynomial evaluation interval is \qtyrange{293.15}{393.15}{K}; this interval does not extend the source ranges of the anchor or reaction correlations.
The CO$_2$ anchor uses the NIST/Chase ideal-gas Shomate relation, whose stated interval begins at \qty{298}{K}; evaluations below that temperature extrapolate it.
The liquid heat-capacity anchors use the MEA and water correlations in \cref{tab:solvent-cp-coefficients} \cite{Hilliard2008}.
Their formation enthalpies at \qty{298.15}{K} are $-393510$, $-507500$, and $-285830$ \unit{J\,mol^{-1}} for gaseous CO$_2$, liquid MEA, and liquid water, respectively, using the NIST/CODATA and Baroody--Carpenter reference records.

\begin{table}[pos=htbp]
\centering\small
\caption{Reference enthalpies at \qty{353.15}{K} for the selected caloric reconstruction. The reference values are in \unit{J\,mol^{-1}}; the functions and coefficients are defined by \cref{eq:caloric-polynomial,tab:caloric-polynomial-coefficients}.}
\label{tab:species-reference-calorics}
\begin{tabular}{@{}llr@{}}
\toprule
Species & Construction & $h_i^\circ(T_0)$\\
\midrule
CO$_2$ & Ideal-gas anchor & -391399.851722\\
MEA & Liquid anchor & -439845.547192\\
H$_2$O & Liquid anchor & -239696.855347\\
MEAH$^+$ & Reaction reconstruction & -1241508.879099\\
MEACOO$^-$ & Reaction reconstruction & 576475.883657\\
HCO$_3^-$ & Reaction reconstruction & 836014.516079\\
CO$_3^{2-}$ & Reaction reconstruction & 3315335.032953\\
H$_3$O$^+$ & Ionic reference convention & -239696.855347\\
OH$^-$ & Reaction reconstruction & 1320557.929724\\
\bottomrule
\end{tabular}
\end{table}

\begin{table}[pos=htbp]
\centering\small
\caption{Coefficients of \cref{eq:caloric-polynomial}, in \unit{J\,mol^{-1}\,K^{-1}}, for dimensionless $\theta$. Ionic reference coefficients follow the reaction and reference constraints; they are not isolated-ion measurements.}
\label{tab:caloric-polynomial-coefficients}
\begin{tabular}{@{}lrrrr@{}}
\toprule
Species & $b_0$ & $b_1$ & $b_2$ & $b_3$\\
\midrule
CO$_2$ & \num{39.5392771} & \num{4.05693353} & \num{-0.516249792} & \num{0.107458931}\\
MEA & \num{108.889073} & \num{17.4811347} & \num{1.26159278} & \num{0.553701791}\\
H$_2$O & \num{43.3004643} & \num{-2.41226427} & \num{0.489801326} & \num{-0.450742496}\\
MEAH$^+$ & \num{111.189062} & \num{36.2721195} & \num{87.2871286} & \num{50.3179557}\\
MEACOO$^-$ & \num{66.5960204} & \num{-70.6414207} & \num{-86.1662791} & \num{-146.388923}\\
HCO$_3^-$ & \num{-65.7974712} & \num{-61.9915653} & \num{-150.449984} & \num{-130.587733}\\
CO$_3^{2-}$ & \num{-149.474152} & \num{-167.136065} & \num{-419.205957} & \num{-360.490126}\\
H$_3$O$^+$ & \num{43.3004643} & \num{-2.41226427} & \num{0.489801326} & \num{-0.450742498}\\
OH$^-$ & \num{-8.08765516} & \num{-70.259169} & \num{-138.651513} & \num{-133.974491}\\
\bottomrule
\end{tabular}
\par\medskip
\begin{tabular}{@{}lrrrr@{}}
\toprule
Species & $b_4$ & $b_5$ & $b_6$ & $b_7$\\
\midrule
CO$_2$ & \num{-0.0352248499} & \num{0.0118568075} & \num{-0.00399193027} & \num{0.00174613689}\\
MEA & \num{-0.475104788} & \num{-0.0589415555} & \num{0.0285678364} & \num{-0.02012004}\\
H$_2$O & \num{0.599298998} & \num{-0.204385847} & \num{0.0647883985} & \num{-0.0243611285}\\
MEAH$^+$ & \num{14.4869303} & \num{-5.2510749} & \num{-8.7129525} & \num{-1.46905663}\\
MEACOO$^-$ & \num{-19.7536858} & \num{14.8060513} & \num{11.8931111} & \num{8.07027862}\\
HCO$_3^-$ & \num{-26.5920306} & \num{13.010693} & \num{15.1457418} & \num{6.14537206}\\
CO$_3^{2-}$ & \num{-76.4595145} & \num{37.5688003} & \num{42.9595557} & \num{16.8471102}\\
H$_3$O$^+$ & \num{0.599298998} & \num{-0.204385835} & \num{0.0647883995} & \num{-0.0243611511}\\
OH$^-$ & \num{-24.3559863} & \num{12.6243717} & \num{15.2487565} & \num{5.94574174}\\
\bottomrule
\end{tabular}
\end{table}
% CHECKLIST-END: method-12

The thermal-conductivity correlation has the form
\begin{align}
k_{t,i}=\frac{A_iT^{B_i}}{1+C_i/T+D_i/T^2},
\end{align}
with $k_{t,i}$ in \unit{W\,m^{-1}\,K^{-1}} and coefficients appropriate to the species and phase.
The heat-transfer relation uses the mixture conductivity consistently with the selected property description.
For the gas mixture, $k_{t,v}=\sum_i\frac{y_i k_{t,i}}{\sum_j y_j\Phi_{ij}}$, using the same viscosity-based $\Phi_{ij}$ as the gas-viscosity relation.

\begin{table}[pos=htbp]
\centering\small
\caption{Gas thermal-conductivity coefficients. With $T$ in kelvin, the stated correlation returns \unit{W\,m^{-1}\,K^{-1}}.}
\label{tab:gas-conductivity-coefficients}
\begin{tabular}{@{}lrrrr@{}}
\toprule
Species & $A$ & $B$ & $C$ & $D$\\
\midrule
CO$_2$ & 3.69 & $-0.3838$ & 964 & $1.86\times10^6$\\
Water & $6.204\times10^{-6}$ & 1.3973 & 0 & 0\\
N$_2$ & $3.31\times10^{-4}$ & 0.7722 & 16.323 & 373.72\\
O$_2$ & $4.50\times10^{-4}$ & 0.7456 & 56.699 & 0\\
\bottomrule
\end{tabular}
\end{table}

\subsection{Nine-species thermodynamic and mobility parameters}
The thermodynamic parameters and mobility inputs are reported separately because they determine different parts of the flux law.
The selected nine-species parameterization is summarized below; source precision is retained in the numerical inputs.
The reaction-constant coefficients refer to \cref{eq:reaction-constants}.
For the water segment diameter,
\begin{align}
\sigma_W(T)=2.7927+10.11\exp(-0.01775T)-1.417\exp(-0.01146T)
\quad[\text{\AA}].
\end{align}
% Source: typed MEA_reactive_epcsaft_2026_09_03/parameters.json, SHA-256 a9186c93759f2e2c02a6c913350ad06a244fff3f82503820c9962b3df8dd40d9; unchanged EOS pure/pair records.

% P:parameters-01
\Cref{tab:reactive-pure-parameters} summarizes the selected segment and ionic diameter parameters for the reactive liquid.
These values define the nine-species pressure/speciation parameterization used to specify the reactive liquid.
The complete typed parameter document also retains association-site topology, pair parameters, correlations, source locators, and qualification ranges in the accompanying computational data.
The water segment diameter uses the temperature correlation stated above.
The selected coefficient set assigns a solvation factor of 1.5 to water and 1 to every other species.

\begin{table}[pos=htbp]
\centering
\caption{Selected nine-species thermodynamic parameters. Diameters are in \unit{\angstrom}; $\epsilon/k$ is in \unit{K}. Packing and Debye--H\"uckel diameters coincide in this parameter set. Displayed values are rounded; calculations use the retained document precision.}
\label{tab:reactive-pure-parameters}
\small
\begin{tabular}{@{}l S[table-format=1.4] c S[table-format=3.3] c c@{}}
\toprule
Species & {$m$} & {$\sigma$} & {$\epsilon/k$} & {$d_{\mathrm{Born}}$} & {$d_{\mathrm{pack}}=d_{\mathrm{DH}}$} \\
\midrule
\ce{CO2} & 2.0729 & 2.7852 & 173.440 & -- & -- \\
\ce{MEA} & 3.0353 & 3.0435 & 277.174 & -- & -- \\
\ce{H2O} & 1.2047 & $\sigma(T)$ & 353.950 & -- & -- \\
\ce{MEAH+} & 1.0000 & 3.4851 & 232.687 & 3.5332 & 3.0669 \\
\ce{MEACOO-} & 1.0000 & 3.5354 & 453.265 & 3.5411 & 3.1112 \\
\ce{HCO3-} & 1.0000 & 2.9296 & 70.000 & 3.0000 & 2.5780 \\
\ce{CO3^2-} & 1.0000 & 2.4422 & 249.260 & 3.0000 & 2.1491 \\
\ce{H3O+} & 1.0000 & 3.4654 & 500.000 & 1.2180 & 3.0496 \\
\ce{OH-} & 1.0000 & 2.0177 & 650.000 & 3.0811 & 1.7756 \\
\bottomrule
\end{tabular}
\end{table}

\begin{table}[pos=htbp]
\centering
\caption{Species diffusivity estimates used in the film mobility, in \unit{m^2\,s^{-1}}. The CO$_2$ expression uses $T$ in kelvin and $R$ in \unit{J\,mol^{-1}\,K^{-1}}. The values specify the species-diffusivity approximation in the mobility law.}
\label{tab:film-diffusivities}
\begin{tabular}{@{}ll@{}}
\toprule
Species & $D_i$ \\
\midrule
\ce{CO2} & $0.5(6.28\times10^{-7})\exp[-15230/(RT)]$ \\
\ce{MEA}, \ce{H2O} & $8.8\times10^{-10}$ \\
\ce{MEAH+} & $8.4\times10^{-10}$ \\
\ce{MEACOO-}, \ce{HCO3-}, \ce{CO3^2-} & $6.8\times10^{-10}$ \\
\ce{H3O+}, \ce{OH-} & $\sqrt{3.4\times10^{-18}}$ \\
\bottomrule
\end{tabular}
\end{table}



% CHECKLIST-BEGIN: method-14
\Cref{tab:film-diffusivities} specifies the central species-diffusivity estimates.
The CO$_2$ expression applies a factor of 0.5 to an aqueous-MEA Arrhenius correlation to represent the selected loading correction.
The MEA and carbamate values provide the molecular and carbamate transport anchors; the protonated-MEA value uses the unresolved MEA/MEAH$^+$ diffusion estimate.
Water is assigned the MEA value, and bicarbonate and carbonate are assigned the carbamate value.
Hydronium and hydroxide use the geometric midpoint of \qty{3.4e-10}{m^2\,s^{-1}} and \qty{1.0e-8}{m^2\,s^{-1}}.
These assignments define a species-diffusivity approximation, rather than measured multicomponent pair diffusivities.
The pair values are formed by the harmonic mean in \cref{eq:film-projection}; the projection then imposes the specified flux constraints.
The thickness diffusivity $D_C^{\mathrm{column}}$ instead uses the common-property liquid CO$_2$ correlation at the bulk state, so changing a species mobility estimate does not implicitly redefine the film thickness.
% CHECKLIST-REVIEW: Add the Polat, Melnikov, and Jerng primary citations with their concentration/temperature ranges; verify final species assignments and sources for the water and ionic estimates against the adopted transport record.
% CHECKLIST-END: method-14

The reaction R5 source correlation spans \qtyrange{273.15}{323.15}{K}; using it above \qty{323.15}{K} is a temperature extrapolation of that relation.
The temperature range of a combined thermodynamic parameterization and the source range of each individual correlation are distinct.
The mobility inputs likewise retain their measurement or estimation basis when interpreting sensitivity to species transport.

